\section{Implementasi Algoritma}

Bagian ini menjelaskan implementasi algoritma pencocokan string yang
digunakan pada Judol Detector. Fokus penjelasan adalah modul di folder
\texttt{src/algorithms} yang bersifat mandiri, tidak bergantung pada API
browser, dan mengembalikan struktur hasil yang konsisten.

Seluruh fungsi pencocokan memiliki keluaran bertipe
\texttt{AlgorithmResult} yang berisi daftar hasil \texttt{DetectionMatch},
jumlah kecocokan, waktu eksekusi, dan (bila relevan) jumlah perbandingan
karakter. Struktur data ini didefinisikan di \texttt{types.ts} agar hasil
setiap algoritma dapat dibandingkan dan digabungkan secara seragam.

\subsection{Struktur Data Hasil Deteksi}

Struktur utama yang digunakan oleh semua algoritma adalah sebagai berikut:

\begin{itemize}
	\item \texttt{MatchAlgorithm} mendefinisikan nama algoritma
	(\texttt{KMP}, \texttt{Boyer-Moore}, \texttt{RegEx},
	\texttt{Weighted-Levenshtein}, \texttt{Aho-Corasick}, \texttt{Rabin-Karp}).
	\item \texttt{DetectionMatch} menyimpan \texttt{keyword}, \texttt{matchedText},
	indeks awal dan akhir, serta informasi tambahan seperti \texttt{score}
	atau \texttt{comparisons} bila tersedia.
	\item \texttt{AlgorithmResult} memuat algoritma asal, daftar hasil,
	jumlah hasil, waktu eksekusi, dan total perbandingan.
\end{itemize}

Dengan struktur ini, modul pemindaian dapat menggabungkan hasil tanpa perlu
mengetahui detail internal setiap algoritma.

\subsection{Knuth-Morris-Pratt (\texttt{kmp.ts})}

Implementasi KMP dibangun melalui dua langkah utama: pembuatan
\textit{failure table} dan proses pemindaian.

\begin{itemize}
	\item \textbf{Failure table.} Fungsi \texttt{buildFailureTable} menghitung
	panjang prefiks terpanjang yang juga merupakan sufiks untuk setiap posisi
	pada pola. Tabel ini digunakan untuk menentukan posisi geser ketika
	terjadi ketidakcocokan, sehingga perbandingan karakter tidak mengulang
	dari awal.

	\item \textbf{Pemindaian teks.} Fungsi \texttt{searchKmp} menelusuri teks
	dan pola dengan dua indeks. Setiap kali pola cocok penuh, posisi awal
	dicatat lalu indeks pola digeser menggunakan failure table.
\end{itemize}

Fungsi publik \texttt{findKmpMatches} melakukan normalisasi sederhana dengan
\texttt{toLowerCase()}, lalu memindai setiap keyword pada teks. Hasil disusun
ke dalam daftar \texttt{DetectionMatch} beserta jumlah perbandingan karakter.

\subsection{Boyer-Moore (\texttt{boyerMoore.ts})}

Algoritma Boyer-Moore menggunakan \textit{bad character heuristic}.
Implementasi ini mencakup dua komponen:

\begin{itemize}
	\item \textbf{Bad character table.} Fungsi \texttt{buildBadCharTable}
	menyimpan posisi kemunculan terakhir setiap karakter pada pola.

	\item \textbf{Pemindaian.} Fungsi \texttt{bmScan} membandingkan karakter
	dari kanan ke kiri. Saat mismatch terjadi, pergeseran dihitung dari
	selisih posisi karakter mismatch di pola dan kemunculan terakhirnya pada
	tabel. Jika karakter tidak ada pada pola, pergeseran dilakukan sejauh
	mungkin.
\end{itemize}

Sebelum pemindaian, teks dan pola dinormalisasi ke bentuk \texttt{NFKC} dan
huruf kecil agar konsisten terhadap variasi Unicode. Hasil akhir disusun
ke \texttt{AlgorithmResult} lengkap dengan jumlah perbandingan karakter.

\subsection{RegEx Matching (\texttt{regexMatcher.ts})}

Pola RegEx yang digunakan adalah:\\
\texttt{(?<![\textbackslash{}p\{L\}\textbackslash{}p\{N\}\_])[\textbackslash{}p\{L\}]\{2,\}[0-9]\{2,3\}(?![\textbackslash{}p\{L\}\textbackslash{}p\{N\}\_])}.
Pola ini mendeteksi token huruf minimal dua karakter yang diikuti dua atau
tiga digit, dengan batas kiri dan kanan agar tidak berada di tengah token
lain.

Fungsi \texttt{findJudolPatternMatches} memanggil \texttt{matchAll} untuk
menemukan semua kemunculan dan menyusun hasil ke \texttt{DetectionMatch}.
Karena RegEx tidak melakukan perbandingan karakter eksplisit, informasi
\texttt{comparisons} tidak diisi.

\subsection{Weighted Levenshtein Distance (\texttt{weightedLevenshtein.ts})}

Implementasi Weighted Levenshtein dibangun dengan tabel dinamis dua dimensi.
Biaya operasi edit menggunakan parameter berikut:

\begin{itemize}
	\item biaya substitusi karakter sama: 0,
	\item biaya substitusi karakter mirip: 0.25 (default),
	\item biaya substitusi karakter berbeda: 1,
	\item biaya insert dan delete: 1.
\end{itemize}

Karakter mirip visual dikelompokkan pada \texttt{VISUAL\_GROUPS}
misalnya \texttt{o/0}, \texttt{a/4}, \texttt{e/3}, \texttt{s/5},
dan variasi Unicode terkait. Fungsi \texttt{getVisualSubstitutionCost}
memutuskan biaya substitusi berdasarkan kelompok tersebut.

Proses pencocokan tidak dilakukan pada seluruh teks, tetapi pada kandidat
token. Token diperoleh melalui pola \texttt{TOKEN\_PATTERN} yang menangkap
huruf, angka, garis bawah, dan simbol tertentu. Untuk keyword multi-kata,
teks dibentuk menjadi jendela token dengan panjang yang sama. Setiap kandidat
yang panjangnya terlalu jauh dari keyword disaring terlebih dahulu untuk
mengurangi biaya komputasi.

Fungsi \texttt{findWeightedLevenshteinMatches} menghitung nilai similarity
berdasarkan $1 - \frac{distance}{maxLength}$ dan membandingkannya dengan
ambang default 0.78. Kecocokan eksak dapat diikutsertakan atau diabaikan
melalui opsi \texttt{includeExact}.

\subsection{Optical Character Recognition (OCR)}

Fitur OCR digunakan untuk mendeteksi teks pada gambar ketika konten judol
tidak muncul sebagai node DOM. Prosesnya dimulai dari pemilihan gambar yang
terlihat pada viewport, lalu mengambil teks konteks dari atribut \texttt{alt},
\texttt{aria-label}, nama file, dan container terdekat. Jika konteks belum
memberikan hasil, sistem menjalankan OCR menggunakan Tesseract.js untuk
mengekstrak teks dari gambar. Teks hasil OCR kemudian diproses oleh pipeline
yang sama (RegEx, exact matching, dan Weighted Levenshtein) agar konsisten
dengan deteksi pada teks DOM.

\subsection{Aho-Corasick (\texttt{ahoCorasick.ts})}

Implementasi Aho-Corasick dibangun dari dua tahap: pembuatan trie dan
pembuatan \textit{failure links}.

\begin{itemize}
	\item \textbf{Trie keyword.} Fungsi \texttt{buildTrie} membangun struktur
	trie dari seluruh keyword, dan setiap node menyimpan daftar keyword yang
	berakhir di node tersebut.

	\item \textbf{Failure links.} Fungsi \texttt{buildFailureLinks} menggunakan
	BFS untuk menambahkan tautan kegagalan yang mengarahkan ke prefiks terpanjang
	yang juga merupakan sufiks.
\end{itemize}

Fungsi \texttt{findAhoCorasickMatches} menelusuri teks satu kali dan, ketika
mencapai node dengan \texttt{outputs}, mencatat seluruh keyword yang cocok.
Algoritma ini cocok untuk pencarian banyak keyword sekaligus dengan biaya
linear terhadap panjang teks dan jumlah hasil.

\subsection{Rabin-Karp (\texttt{rabinKarp.ts})}

Rabin-Karp pada implementasi ini menggunakan \textit{rolling hash} dengan
basis 31 dan modulo \texttt{1000000007}. Prosesnya meliputi:

\begin{itemize}
	\item menghitung hash pola dan hash jendela awal,
	\item menggeser jendela dengan menghapus karakter kiri dan menambahkan
	karakter kanan (rolling),
	\item melakukan verifikasi karakter demi karakter jika nilai hash sama
	untuk mencegah \textit{hash collision}.
\end{itemize}

Fungsi \texttt{findRabinKarpMatches} melakukan normalisasi \texttt{NFKC} dan
huruf kecil sebelum pemindaian. Jumlah perbandingan dihitung dari jumlah
pemeriksaan hash dan verifikasi karakter.

\subsection{Catatan Implementasi Umum}

Beberapa keputusan implementasi dilakukan agar hasil algoritma konsisten dan
siap digabungkan oleh pipeline pemindaian:

\begin{itemize}
	\item Normalisasi huruf kecil dilakukan pada KMP, Boyer-Moore,
	Aho-Corasick, dan Rabin-Karp untuk membuat pencocokan tidak sensitif
	terhadap kapitalisasi.
	\item Normalisasi Unicode \texttt{NFKC} digunakan pada Boyer-Moore dan
	Rabin-Karp untuk menyamakan representasi karakter yang tampak sama.
	\item Semua hasil selalu menyimpan \texttt{startIndex} dan \texttt{endIndex}
	pada teks asli agar proses highlight dapat bekerja tanpa kehilangan
	konteks.
	\item Pengukuran waktu eksekusi dilakukan dengan \texttt{performance.now()}
	agar presisi cukup untuk membandingkan algoritma pada skenario nyata.
\end{itemize}

Bagian ini menjadi fondasi implementasi pipeline deteksi pada content script
karena semua algoritma menyediakan format hasil yang seragam dan dapat
digabungkan melalui mekanisme deduplikasi dan statistik.

\section{Spesifikasi Teknis Program}

Bagian ini merangkum struktur data, fungsi utama, dan prosedur yang dibangun
pada sistem.

\subsection{Struktur Data}

\begin{table}[H]
\centering
\caption{Ringkasan struktur data utama}
\label{tab:struktur-data}
\begin{tabular}{|p{4cm}|p{9cm}|}
\hline
\textbf{Struktur} & \textbf{Deskripsi} \\
\hline
\texttt{MatchAlgorithm} & Enumerasi nama algoritma yang digunakan. \\
\hline
\texttt{DetectionMatch} & Hasil satu kecocokan: keyword, teks terdeteksi, indeks, dan metadata. \\
\hline
\texttt{AlgorithmResult} & Ringkasan hasil satu algoritma: daftar match, count, waktu eksekusi. \\
\hline
\end{tabular}
\end{table}

\subsection{Fungsi dan Prosedur Utama}

Fungsi dan prosedur utama yang dibangun meliputi:

\begin{itemize}
	\item \texttt{findKmpMatches}: pencocokan eksak menggunakan KMP.
	\item \texttt{buildFailureTable}: membangun failure table untuk KMP.
	\item \texttt{findBoyerMooreMatches}: pencocokan eksak menggunakan Boyer-Moore.
	\item \texttt{buildBadCharTable}: tabel kemunculan terakhir karakter pada Boyer-Moore.
	\item \texttt{findJudolPatternMatches}: deteksi pola RegEx \texttt{<kata><angka>}.
	\item \texttt{findWeightedLevenshteinMatches}: pencocokan fuzzy berbobot untuk algoritma Weighted Levenshtein.
	\item \texttt{calculateWeightedLevenshtein}: menghitung jarak berbobot dan similarity untuk algoritma Weighted Levenshtein.
	\item \texttt{getVisualSubstitutionCost}: biaya substitusi untuk karakter mirip visual untuk Weighted Levenshtein.
	\item \texttt{findAhoCorasickMatches}: pencocokan multi-keyword dengan trie untuk algoritma Aho-Corasick.
	\item \texttt{buildTrie} dan \texttt{buildFailureLinks}: konstruksi trie dan failure links untuk algoritma Aho-Corasick.
	\item \texttt{findRabinKarpMatches}: pencocokan berbasis rolling hash untuk algoritma Rabin-Karp.
	\item \texttt{rkScan} dan \texttt{computeHash}: pemindaian dan hash window Rabin-Karp.
\end{itemize}

\subsection{Prosedur Utama Pipeline}

\begin{enumerate}
	\item Membaca daftar keyword dari \texttt{keyword.txt}.
	\item Mengambil node teks dari DOM yang terlihat.
	\item Menjalankan pencocokan eksak (KMP, Boyer-Moore, Aho-Corasick, Rabin-Karp).
	\item Menjalankan RegEx untuk pola \texttt{<kata><angka>}.
	\item Menjalankan Weighted Levenshtein pada kandidat token.
	\item Melakukan deduplikasi hasil dan menyiapkan statistik.
	\item Mengirim hasil ke modul highlight dan popup.
\end{enumerate}

\section{Tata Cara Penggunaan Extension}

Bagian ini menjelaskan antarmuka dan fitur yang tersedia bagi pengguna.

\subsection{Memuat Extension}

\begin{enumerate}
	\item Jalankan proses build hingga menghasilkan folder \texttt{dist/}.
	\item Buka \textit{Extensions} pada Chromium dan aktifkan \textit{Developer Mode}.
	\item Pilih \textit{Load unpacked} lalu arahkan ke folder \texttt{dist/}.
\end{enumerate}

\subsection{Antarmuka Popup}

\begin{itemize}
	\item Tombol \textbf{Scan/Rescan} untuk memulai pemindaian.
	\item Toggle \textbf{Blur} untuk menyamarkan konten terdeteksi.
	\item Toggle \textbf{OCR} untuk memindai teks pada gambar (jika tersedia).
	\item \textbf{Filter highlights} untuk memilih algoritma yang ditampilkan secara
	\textit{intersection}.
	\item Ringkasan statistik: total match, per algoritma, dan waktu eksekusi.
\end{itemize}

\subsection{Interaksi pada Halaman}

\begin{itemize}
	\item Elemen terdeteksi diberi highlight.
	\item Hover pada highlight menampilkan tooltip berisi keyword, algoritma,
	jumlah kemunculan, dan waktu eksekusi.
	\item Rescan menghapus highlight lama dan memperbarui statistik.
\end{itemize}

\section{Pengujian}

Bagian ini memuat beberapa kasus uji dengan variasi algoritma dan jenis
halaman. Tabel berikut adalah template yang dapat diisi dengan hasil aktual.

\begin{table}[H]
\centering
\caption{Daftar kasus uji}
\label{tab:kasus-uji}
\begin{tabular}{|p{1.2cm}|p{5.5cm}|p{6cm}|}
\hline
\textbf{ID} & \textbf{Jenis Halaman} & \textbf{Algoritma Utama} \\
\hline
TC-01 & haloowner.pages.dev & KMP dan Aho-Corasick \\
\hline
TC-02 & \"maxwin234\" di search bar & RegEx \\
\hline
TC-03 & heylink.me/88omegabet & Weighted Levenshtein dan RegEx \\
\hline
TC-04 & \"AB1 A9\" di search bar & RegEx \\
\hline
TC-05 & one-vv5378.com/casino?p=3o0j & Boyer-Moore dan Rabin-Karp \\
\hline
TC-06 & tumgik.com/promaxwin555-page & Semua algoritma \\
\hline
\end{tabular}
\end{table}

\subsection{Dokumentasi Kasus Uji}

\begin{figure}[H]
	\centering
	\includegraphics[width=0.92\textwidth]{pengujian/TC1_all.png}
	\caption{TC-01: Hasil pemindaian dengan semua algoritma aktif.}
	\label{fig:tc1-all}
\end{figure}

\begin{figure}[H]
	\centering
	\includegraphics[width=0.92\textwidth]{pengujian/TC1_KMP.png}
	\caption{TC-01: Filter highlight hanya KMP.}
	\label{fig:tc1-kmp}
\end{figure}

\begin{figure}[H]
	\centering
	\includegraphics[width=0.92\textwidth]{pengujian/TC1_AC.png}
	\caption{TC-01: Filter highlight hanya Aho-Corasick.}
	\label{fig:tc1-ac}
\end{figure}

\begin{figure}[H]
	\centering
	\includegraphics[width=0.92\textwidth]{pengujian/TC2_regex.png}
	\caption{TC-02: Deteksi pola \texttt{MAXWIN234} oleh RegEx.}
	\label{fig:tc2}
\end{figure}

\begin{figure}[H]
	\centering
	\includegraphics[width=0.92\textwidth]{pengujian/TC2_keyword.png}
	\caption{TC-02: Hasil perbandingan keyword.}
	\label{fig:tc2}
\end{figure}

\begin{figure}[H]
	\centering
	\includegraphics[width=0.92\textwidth]{pengujian/TC3_WL.png}
	\caption{TC-03: Deteksi fuzzy dengan Weighted Levenshtein}
	\label{fig:tc3-wl}
\end{figure}

\begin{figure}[H]
	\centering
	\includegraphics[width=0.92\textwidth]{pengujian/TC3_regex.png}
	\caption{TC-03: Deteksi dengan RegEx}
	\label{fig:tc3-wl}
\end{figure}

\begin{figure}[H]
	\centering
	\includegraphics[width=0.92\textwidth]{pengujian/TC4.png}
	\caption{TC-04: Contoh false positive pola RegEx pada token umum.}
	\label{fig:tc4}
\end{figure}

\begin{figure}[H]
	\centering
	\includegraphics[width=0.92\textwidth]{pengujian/TC5_all.png}
	\caption{TC-05: Deteksi dengan semua algoritma dalam webpage.}
	\label{fig:tc5-all}
\end{figure}

\begin{figure}[H]
	\centering
	\includegraphics[width=0.92\textwidth]{pengujian/TC5_BM.png}
	\caption{TC-05: Deteksi dengan algoritma Boyer-Moore.}
	\label{fig:tc5-bm}
\end{figure}

\begin{figure}[H]
	\centering
	\includegraphics[width=0.92\textwidth]{pengujian/TC5_RK.png}
	\caption{TC-05: Deteksi dengan semua algoritma dalam webpage.}
	\label{fig:tc5-rk}
\end{figure}

\begin{figure}[H]
	\centering
	\includegraphics[width=0.92\textwidth]{pengujian/TC6_all.png}
	\caption{TC-06: Deteksi dengan semua algoritma dalam webpage.}
	\label{fig:tc6-all}
\end{figure}

\begin{figure}[H]
	\centering
	\includegraphics[width=0.92\textwidth]{pengujian/OCR.png}
	\caption{OCR: Deteksi teks pada gambar hasil pencarian gambar.}
	\label{fig:ocr}
\end{figure}

\subsection{Hasil Pengujian}

\begin{table}[H]
\centering
\caption{Hasil pengujian}
\label{tab:hasil-uji}
\begin{tabular}{|p{1.5cm}|p{4.2cm}|p{3.5cm}|p{2.5cm}|}
\hline
\textbf{ID TC} & \textbf{Jenis Algoritma} & \textbf{Banyak Match} & \textbf{Waktu (ms)} \\
\hline
TC-01 & KMP & 8 & 12.10 \\
\hline
TC-01 & Aho-Corasick & 8 & 22.8 \\
\hline
TC-02 & RegEx & 23 & 7.0 \\
\hline
TC-03 & RegEx & 4 & 8.8 \\
\hline
TC-03 & Weighted Levenshtein & 4 & 7373.0 \\
\hline
TC-04 & RegEx & 14 & 6.5 \\
\hline
TC-05 & Boyer-Moore & 5 & 16.1 \\
\hline
TC-05 & Rabin-Karp & 5 & 23.5 \\
\hline
TC-06 & KMP & 14 & 33.3 \\
\hline
TC-06 & Aho-Corasick & 14 & 52.2 \\
\hline
TC-06 & RegEx & 10 & 5.1 \\
\hline
TC-06 & Boyer-Moore & 15 & 12.9 \\
\hline
TC-06 & Rabin-Karp & 15 & 25.4 \\
\hline
TC-06 & Weighted Levenshtein & 4 & 2122.8 \\
\hline
\end{tabular}
\end{table}

\section{Analisis Hasil Pengujian}

Bagian ini membahas perbandingan performa dan akurasi antar algoritma.

\subsection{Perbandingan Performa}

\begin{itemize}
	\item RegEx konsisten menjadi yang tercepat pada semua kasus uji
	(5.1--8.8 ms), karena hanya melakukan pencocokan pola tunggal pada teks.
	Keunggulan: sangat cepat untuk pola sederhana. Kelemahan: mudah memicu
	false positive tanpa konteks.
	\item Algoritma exact matching (KMP, Boyer-Moore, Aho-Corasick, Rabin-Karp)
	berada pada rentang puluhan milidetik; KMP cenderung lebih cepat daripada
	Aho-Corasick pada kasus dengan keyword terbatas (TC-01), sedangkan
	Aho-Corasick tetap stabil ketika jumlah keyword bertambah (TC-06).
	Keunggulan KMP: stabil dan sederhana untuk satu pola. Kelemahan: perlu
	memindai per keyword. Keunggulan Boyer-Moore: lompatan besar saat mismatch,
	baik untuk pola pendek. Kelemahan: performa bisa turun pada teks repetitif.
	Keunggulan Aho-Corasick: satu lintasan untuk banyak keyword. Kelemahan:
	overhead trie dan failure links. Keunggulan Rabin-Karp: hash cepat untuk
	window. Kelemahan: tetap perlu verifikasi saat hash cocok.
	\item Weighted Levenshtein memiliki waktu paling besar (2.1--7.3 detik)
	karena menghitung jarak berbobot pada banyak kandidat token; hal ini
	menunjukkan kebutuhan pembatasan kandidat agar tetap responsif.
	Keunggulan: menangkap variasi visual yang tidak terdeteksi eksak.
	Kelemahan: biaya komputasi tinggi pada teks panjang.
\end{itemize}

\subsection{Pembahasan Akurasi dan False Positive}

\begin{itemize}
	\item Pencocokan eksak (KMP, Boyer-Moore, Aho-Corasick, Rabin-Karp)
	menunjukkan hasil konsisten pada keyword yang benar, terutama pada TC-01
	(8 match) dan TC-06 (14--15 match).
	\item RegEx efektif menangkap pola \texttt{<kata><angka>} pada TC-02 dan
	TC-03, namun menghasilkan false positive pada TC-04 (contoh token umum
	\texttt{AB1}, \texttt{A9}, \texttt{Rp33}, \texttt{Rp32}). Hal ini wajar
	karena pola RegEx belum mempertimbangkan konteks semantik.
	\item Weighted Levenshtein berhasil menemukan variasi visual (TC-03),
	namun biaya komputasinya tinggi. Threshold default 0.78 cukup sensitif
	untuk variasi karakter, tetapi perlu dikontrol agar tidak memicu terlalu
	banyak kandidat pada halaman besar.
	\item OCR memperluas cakupan deteksi ke konten gambar (lihat Gambar
	\ref{fig:ocr}). Hasil OCR efektif untuk menangkap teks judol yang tidak
	tersedia sebagai DOM, namun akurasinya bergantung pada kualitas gambar
	dan bisa menghasilkan teks parsial; karena itu tetap diperlukan validasi
	lewat pipeline string matching untuk mengurangi false positive.
\end{itemize}

\subsection{Ringkasan Analisis}

Secara keseluruhan, pipeline mampu mendeteksi keyword dan pola judol pada
beragam halaman, dengan RegEx dan algoritma exact matching sebagai komponen
tercepat dan paling stabil. Weighted Levenshtein terbukti penting untuk
variasi visual, tetapi perlu digunakan selektif agar performa tetap baik.
Untuk mengurangi false positive, pola RegEx sebaiknya ditambah filter
berdasarkan prefix umum atau konteks kata kunci sekitar.
Keunggulan utama RegEx adalah kecepatan, sementara kelemahannya ada pada
akurasi konteks. KMP dan Boyer-Moore unggul untuk pencocokan eksak dengan
overhead rendah, namun skalanya terbatas saat keyword banyak. Aho-Corasick
memberi keuntungan ketika keyword banyak, tetapi membutuhkan preprocessing
lebih berat. Rabin-Karp efektif untuk pemeriksaan berbasis hash, tetapi
masih memerlukan verifikasi agar aman dari collision. Weighted Levenshtein
memberi fleksibilitas terhadap variasi visual, tetapi menjadi bottleneck
performa tanpa pembatasan kandidat.
